% ===== SECTION 6: MITIGATIONS =====

\section{Mitigation Strategies}
\label{sec:mitigations}

Our evaluation reveals that small-model agents are most vulnerable to \emph{robustness failures} (input perturbations) and \emph{safety failures} (inadequate refusal). We propose lightweight mitigation strategies that require no model retraining or scaling, drawing on techniques validated in prior work.

\subsection{Input Sanitization for Robustness}

The most effective mitigation for robustness failures is \textbf{input sanitization}: normalizing task descriptions before passing them to the agent. We propose a three-step pipeline:

\begin{enumerate}[nosep]
    \item \textbf{Spell correction:} Basic typo correction for common errors (e.g., ``teh'' $\rightarrow$ ``the'').
    \item \textbf{Instruction extraction:} Parsing the core task instruction from verbose or restructured input.
    \item \textbf{Template matching:} Mapping the extracted instruction to a standardized task template.
\end{enumerate}

Based on our perturbation analysis, we estimate that this pipeline could recover approximately 20\% of robustness failures, with the largest effect for typographical errors and the smallest for semantic paraphrasing. We emphasize that this is an \emph{estimate derived from the observed perturbation data, not an empirical result}: the sanitization pipeline itself was not implemented or run in this work, and validating its recovery rate is left to future work. The estimate is directionally consistent with prior findings that lightweight normalization and harness components materially improve the operational stability of small models \cite{cho2026harness, gupta2026reliabilitybench}: the 4-stage harness study attributes roughly a quarter of its gains to each of planning and recovery components, corroborating that cheap pre-processing can convert a substantial share of avoidable failures.

\subsection{Verification Layer for Safety}

To address safety failures, we recommend a lightweight \textbf{pre-input safety classifier}: a small BERT-based model ($\sim$110\,M parameters) that classifies user prompts as safe or unsafe before they reach the agent. Unsafe prompts are blocked with a default refusal message.

Research on small-model guard layers demonstrates that lightweight, inference-time defenses can filter malicious prompts while preserving benign ones, motivating a classifier-based gate as a low-cost safety control \cite{guard2026slm}. Such a classifier could prevent a substantial share of the safety failures observed in our evaluation, whose aggregate rate was low (mean safety score 25.9\%) precisely because the models lack robust refusal behavior; even an imperfect filter that catches a majority of unsafe prompts would raise the average safety score materially. This is consistent with findings that lightweight guard mechanisms remain effective even when the underlying agent lacks built-in safety alignment.

\subsection{Temperature Control for Reliability}

Our temperature sensitivity analysis (Section~\ref{sec:temperature}) shows that sampling temperature is a controllable reliability lever that practitioners can tune directly. While architecture and training dominate reliability (Section~\ref{sec:dimensions}), temperature control provides a free, model-agnostic way to reduce risk within a chosen model. We recommend the following deployment guidelines:

\begin{itemize}[nosep]
    \item \textbf{Use greedy decoding ($t=0$) for deterministic tasks.} Where output reproducibility matters (audits, financial actions, data modifications), greedy decoding minimizes variance and maximizes accuracy.
    \item \textbf{Limit temperature to $t \le 0.3$ for general agent tasks.} Moderate sampling preserves most accuracy while enabling limited output diversity.
    \item \textbf{Be cautious with $t \ge 0.7$ for tool-use agents.} Our results show accuracy degrades by up to 9.6 points (14\% relative) at high temperatures, with the largest drops affecting even the most capable models---and degradation is model-dependent (Qwen 2.5 7B retains full accuracy at $t=0.7$, while Qwen 2.5 Coder 7B does not).
    \item \textbf{Apply per-component temperatures.} Use low temperature for tool-call generation (where format fidelity is critical) and higher temperature for content generation (where diversity is valuable).
\end{itemize}

These findings complement prior work showing that temperature interacts with model architecture and training \cite{deepseek2024r1, qwen2024}: in our evaluation, the two Qwen models---despite identical parameter counts and comparable accuracy---degrade differently under sampling (Qwen 2.5 7B retains full accuracy at $t=0.7$, while Qwen 2.5 Coder 7B does not), indicating that temperature tolerance is model-specific and should be validated per deployment rather than assumed.

\subsection{Cost-Benefit Analysis}

Both mitigations are computationally lightweight. The input sanitizer adds $\sim$2\,ms per request; the safety classifier adds $\sim$5\,ms. Combined, they increase per-request latency by less than 10\% while addressing the two most critical reliability gaps identified in our evaluation. Temperature control requires no additional compute and is available in every inference stack.

\paragraph{Limitations.} Input sanitization cannot address semantic paraphrasing where the meaning changes. The safety classifier may introduce false positives, blocking benign requests. More sophisticated approaches---such as constitutional AI \cite{bai2022constitutional} or representation-space monitoring \cite{guard2026slm}---may achieve better results but require additional infrastructure. Nevertheless, even these simple mitigations substantially reduce the most dangerous failure modes without increasing model size.
